% il existe plusieures classes de documents
% pour des documents plus longs, vous pouvez utiliser
% book ou report
\documentclass[a4paper]{article} 



% vous pouvez changer les paramètres : voici les options dispoinibles :
% - a4paper
% - fancysections
% - notitlepage ou titlepage
% - onside ou twoside selon si vous voulez l'imprimer en recto-verso ou en recto
% - sectionmark
% - chaptermark (pour les 
% - pagenumber
% - enmanquedinspiration
% en cas de doutes, pas de doutes, la documentation est sur :
%  https://gitlab.binets.fr/typographix/polytechnique/-/blob/master/source/polytechnique.pdf
\usepackage[a4paper,  fancysections,  titlepage]{polytechnique}
\usepackage[french]{babel}
\usepackage[T1]{fontenc}
\usepackage{blindtext}
\usepackage[hidelinks]{hyperref}
\usepackage{amsmath, amssymb}



% nous avons défini deux commandes :
\newcommand{\code}[1]{%
    \mbox{\ttfamily%
        \detokenize{#1}%
    }%
}

\newcommand{\resultat}[1]{%
    \quad \rightsquigarrow \quad #1%
}

\author{
	Côme Rochas, Thomas Gastellu\\ 
	Benjamin Selles, Leonel Gomes Abreu\\
	Théo Clary
}
\date{\today}
\title{Pokemon RL Project}%
\subtitle{EA Reinforcement Learning}%
% pour changer de logo, ajoutez l'image dans un fomat PDF
% ou image en la glissant à droite et remplacez typographix
% par le nom de l'image, si vous ne voulez pas de logo, 
% supprimze la ligne. 
\logo{logo}

\begin{document}
	\maketitle 
	\tableofcontents
	\pagebreak
	\section{Introduction}

In this project, we explore the application of reinforcement learning (RL) techniques to the domain of Pokemon battles. 
Our goal is to develop an RL agent capable of learning effective strategies for battling against a heuristic opponent.

The pokemon battle environment is a turn-based game where two players select actions (such as attacking, switching Pokemon, or using items) to defeat their opponent's team.
The action space is also large, as there are multiple moves and strategies to consider at each turn.
This environment presents three main challenges:
\begin{itemize}
    \item \textbf{High-dimensional state and action spaces}: The large number of possible states and actions makes it difficult for the agent to learn effective policies.
    \item \textbf{Partial observability}: The agent does not have access to the full state of the environment, as it cannot see the opponent's Pokemon or their moves.
    \item \textbf{High stochasticity}: The outcome of actions can be highly variable due to factors such as move accuracy, critical hits, and random effects, making it challenging for the agent to learn consistent strategies.
\end{itemize}

Previous approaches range from simple scripted baselines to modern reinforcement learning and search-based agents. 
Common baselines include hand-crafted heuristics such as always picking the highest base-power move or switching when at a clear type disadvantage, as well as purely random policies used for benchmarking. 
Classical methods like tabular Q-learning and softmax exploration have been applied in simplified environments \cite{kalose2018optimal}, showing consistent gains over random agents while remaining limited by state-space size. 
More recent work uses deep RL (DQN, A3C, PPO) on top of Pokémon Showdown via dedicated environments, often finding policy-gradient methods like PPO particularly effective in stochastic, high-dimensional settings. 
Beyond RL, some systems combine minimax-style lookahead with learned value functions or leverage large language models \cite{jain2025large} and offline sequence models trained on battle logs, which excel at capturing long-horizon strategy and opponent modeling but trade off interpretability and computational cost.

We have implemented two PPO-based agents: a simple policy network and a more complex actor-critic architecture.
The report is structured as follows: we first describe the environment \ref{section:environment}., then we detail the models we implemented \ref{section:models}., followed by the training setup \ref{section:training}. and finally we present our experimental results \ref{section:results}.
	\pagebreak
	\section{Environment}
\label{section:environment}

The environment is implemented with \texttt{poke-env}, connected to a local Pokemon Showdown server.
In our code, \texttt{PokeRLEnv} extends \texttt{SinglesEnv} (multi-agent), and is wrapped by
\texttt{MaskableSingleAgentWrapper} to expose a standard single-agent \texttt{gymnasium.Env}
interface for Stable-Baselines3.

At each decision step, \texttt{poke-env} provides a \texttt{Battle} object (state snapshot of the
current turn). Our observation pipeline is:
\[
\texttt{Battle state} \;\longrightarrow\; \texttt{embed\_battle(battle)} \;\longrightarrow\; \mathbf{s}_t \in \mathbb{R}^{490}.
\]

The \texttt{Battle} object is not a flat vector: it is a structured Python object with many fields.
Our encoder reads only the fields needed for policy learning:
\begin{itemize}
	\item active Pokemon on both sides (HP fraction, types, boosts, status),
	\item known moves of the active Pokemon,
	\item revealed team members on both sides (bench information),
	\item side conditions (entry hazards and screens),
	\item global effects (weather and field effects),
	\item availability flags for battle gimmicks (Dynamax, Mega, Z-Move, Tera).
\end{itemize}

This is still a partially observable setting: opponent information is only available when revealed
during battle (e.g., observed Pokemon and seen moves).

The state vector is built by concatenating feature blocks in a fixed order:
\[
\mathbf{s}_t = [\mathbf{s}^{ally\_active},\; \mathbf{s}^{ally\_moves},\; \mathbf{s}^{opp\_active},\; \mathbf{s}^{opp\_moves},\; \mathbf{s}^{ally\_bench},\; \mathbf{s}^{opp\_bench},\; \mathbf{s}^{side},\; \mathbf{s}^{weather},\; \mathbf{s}^{field},\; \mathbf{s}^{gimmick}].
\]

\begin{itemize}
	\item \textbf{Ally active Pokemon (33)}: HP fraction (1), type multi-hot over 18 types,
	7 normalized stat boosts (\(-6\) to \(+6\) scaled to \([-1,1]\)), and status one-hot (7).

	\item \textbf{Ally active moves (92)}: 4 fixed slots, each encoded with 23 values:
	move type one-hot (18), base power/250, accuracy/100, category bits (physical/special),
	and normalized PP \(\texttt{current\_pp}/\texttt{max\_pp}\).

	\item \textbf{Opponent active Pokemon (33)}: same encoding as ally active Pokemon.

	\item \textbf{Known opponent moves (220)}: up to 10 known moves, 22 values each
	(same as above \emph{without} PP, because opponent PP is not exposed by the server).
	Unknown slots are zero-padded.

	\item \textbf{Ally bench (40)} and \textbf{opponent bench (40)}:
	5 slots per side, each slot has HP fraction (1) + status one-hot (7), for 8 values per slot.
	Missing/unrevealed slots are zero-padded.

	\item \textbf{Side conditions (14)}:
	for each side, 4 hazard values (Stealth Rock, Spikes layers, Toxic Spikes layers, Sticky Web)
	and 3 binary screen flags (Reflect, Light Screen, Aurora Veil).

	\item \textbf{Weather (8)}: one-hot over the supported weather states.

	\item \textbf{Field effects (6)}: one-hot over terrain/global field effects
	(Electric/Grassy/Misty/Psychic Terrain, Trick Room, Gravity).

	\item \textbf{Gimmick availability (4)}:
	\([\texttt{can\_dynamax},\texttt{can\_mega\_evolve},\texttt{can\_tera},\texttt{can\_z\_move}]\).
\end{itemize}

Therefore, the final observation size is:
\[
33 + 92 + 33 + 220 + 40 + 40 + 14 + 8 + 6 + 4 = 490.
\]

All features are cast to \texttt{float32}. Continuous values are normalized (HP, boosts,
power, accuracy, PP), categorical values are encoded as one-hot/multi-hot vectors, and variable-length
information (moves, bench, revealed opponent data) is converted to fixed-size blocks by zero-padding.

	\pagebreak
	\section{Models}
\label{section:models}

Our core model is \textbf{masked PPO}, implemented with \texttt{MaskablePPO} from
\texttt{sb3-contrib} (Stable-Baselines3 ecosystem), with \texttt{RecurrentPPO} available as an optional variant.
Using an established SB3 implementation was an intentional choice: it gives us a reliable optimization pipeline,
well-tested training code, and straightforward integration with our Gymnasium-compatible environment.

This choice is also motivated by the structure of Pokemon battles: the state is high-dimensional,
transitions are stochastic, and the action space is discrete with many context-dependent invalid actions.
In this regime, an on-policy actor-critic approach with conservative policy updates is generally more stable
than value-based baselines, and action masking improves sample efficiency by restricting exploration to legal
decisions only.

Operationally, at each turn the policy receives the encoded state vector \(\mathbf{s}_t\) from
Section~\ref{section:environment}, then outputs a categorical distribution over actions.
The environment mask is applied so only legal actions remain selectable; the policy samples (or chooses)
an action from that masked distribution. In parallel, the critic estimates \(V_\phi(\mathbf{s}_t)\), used
to compute an advantage estimate \(\hat{A}_t\) for policy optimization.
This actor-critic decomposition is useful here because it separates action selection from state evaluation:
the actor improves tactical choice quality, while the critic provides a lower-variance learning signal.

PPO updates the policy using the probability ratio
\(r_t(\theta)=\pi_\theta(a_t\mid s_t)/\pi_{\theta_{\text{old}}}(a_t\mid s_t)\), with clipped objective:
\[
L^{\text{clip}}(\theta)=\mathbb{E}_t\left[\min\left(r_t(\theta)\hat{A}_t,\;\mathrm{clip}(r_t(\theta),1-\epsilon,1+\epsilon)\hat{A}_t\right)\right].
\]
The clipping term (cut-off) prevents overly large policy updates in one step.
The final training loss also includes a value-function term and an entropy bonus to keep exploration.

In practice, the clipped ratio term is critical: without it, policy updates can become too aggressive,
especially in noisy environments where short-term returns fluctuate strongly from one rollout to another.
The clipped objective keeps updates within a trust region around the previous policy, which tends to produce
more stable training dynamics.

For completeness, we also implemented additional approaches in the same training framework
(notably \texttt{DQN}, and the recurrent PPO variant). However, these alternatives were not evaluated as
extensively as the default masked PPO setup, so our main analysis and conclusions focus on the SB3 masked PPO model.

In short, the policy learns to assign probability mass to strong legal actions while remaining conservative
between updates, which is particularly important in noisy, partially observable battles.
Implementation details (network architecture and hyperparameters) are reported in Appendix~\ref{section:appendix}.
	\pagebreak
	\section{Training Setup}
\label{section:training}
	\pagebreak
	\section{Experimental Results}
\label{section:results}
	\pagebreak
	\section{Appendix}
\label{section:appendix}

\subsection*{Model architecture}

The PPO variants are configured with an MLP policy:
\texttt{MlpPolicy} with \texttt{net\_arch = [256, 256]}.
The policy and value heads both operate on the 490-dimensional observation vector.
The default algorithm is \texttt{MaskablePPO}; \texttt{RecurrentPPO} can be selected as an alternative
through Hydra (\texttt{algorithm=recurrentppo}).

\subsection*{PPO hyperparameters used}

The current parameters from \texttt{pokerl/configs/models/ppo.yaml} are:
\begin{itemize}
    \item learning rate: \(1\times10^{-4}\)
    \item rollout length \texttt{n\_steps}: 1024
    \item mini-batch size: 128
    \item optimization epochs per rollout: 5
    \item discount factor \(\gamma\): 0.99
    \item GAE parameter \(\lambda\): 0.95
    \item clipping range \(\epsilon\): 0.2
    \item entropy coefficient: 0.02
    \item value-function coefficient: 0.5
    \item max gradient norm: 0.5
\end{itemize}

\subsection*{Training defaults}

Default training configuration (\texttt{pokerl/configs/training/default.yaml}):
\(10^6\) timesteps, heuristic opponent for training, evaluation every 10,000 steps,
and 50 evaluation episodes per checkpoint.

	\pagebreak
	\begin{thebibliography}{}
		\bibitem{kalose2018optimal}{Kalose, Akshay and Kaya, Kris and Kim, Alvin. 2018. Optimal battle strategy in pokemon using reinforcement learning. \emph{Web: https://web. stanford. edu/class/aa228/reports/2018/final151. pdf}.}
		\bibitem{jain2025large}{Jain, Daksh and Jain, Aarya and Desai, Ashutosh and Verma, Avyakt and Bhanuka, Ishan and Narang, Pratik and Kumar, Dhruv. 2025. Large Language Models as Pok$\backslash$'emon Battle Agents: Strategic Play and Content Generation. \emph{arXiv preprint arXiv:2512.17308}.}
	\end{thebibliography}
\end{document}